% HISTORICAL EXTENSION -- NOT A CURRENT COMPACT-SUBMISSION AUTHORITY.
% Orbit-selection sufficiency needs sign consistency on the whole evidence
% fibre; the excluded H/ratio-rate extension is not fully verified. See
% audits/THEOREM_NOVELTY_BIBLIOGRAPHY_REVIEW.md, Section 10, and
% formal/KBound/Probability/ChannelCounterexample.lean.
% ============================================================================
%  THE KNOWABILITY DICHOTOMY for label-free sign-of-benefit, and its causal map.
%  Companion to sections/knowability_capacity_general.tex (the layered general
%  result whose Conjecture C.* this section resolves as far as it honestly can).
%
%  Status labels used throughout:
%     (PROVEN)            -- complete elementary proof, no extra hypothesis.
%     (PROVEN, COND.)     -- proven UNDER an explicitly stated condition.
%     (NUMERICAL)         -- validated by Monte-Carlo only (stated as such).
%     (CONJECTURE/OPEN)   -- precise statement + the exact missing piece.
%
%  All numbers validated by
%     experiments/kbound/theory_validation/val_knowability_dichotomy.py
%     -> results_knowability_dichotomy.json   (fixed seed 20260619).
%
%  HONEST HEADLINE (see Remark rmk:dich-scope):
%   Q1.  A SINGLE structural property Phi (benefit-sign factorization through the
%        observable law) gives a COMPLETE, DISTRIBUTION-FREE dichotomy for the
%        EXISTENCE of a label-free knowability capacity, with the universal
%        threshold tau=1. PROVEN both directions. The only residue is a
%        COMPUTABILITY clause on the frontier margin (automatic when the shift
%        family is identifiable from Q_X); it is NOT a separate regime and it is
%        NOT the monotonicity/MLR condition of the prior section -- those were
%        sufficient, not necessary, and we prove tau=1 survives multi-flip,
%        non-monotone frontiers. So Q1 is a genuine single-property dichotomy at
%        the level of EXISTENCE; it becomes a short 2-level refinement only if one
%        insists the capacity be a CLOSED-FORM computable margin.
%   Q2.  Phi maps to causal structure UNDER A STATED GENERICITY CONDITION:
%        in an SCM where the shift is an intervention, Phi holds for the
%        anticausal-with-invariant-mechanism direction and FAILS generically for
%        the causal-with-shifted-mechanism direction. PROVEN one inclusion
%        cleanly; the biconditional holds only modulo a non-degeneracy
%        (faithfulness-type) genericity assumption, with an explicit
%        counterexample showing the unconditional biconditional is FALSE.
% ============================================================================
\section{The knowability dichotomy and its causal correspondence}
\label{sec:knowdich}

% local macro so this section compiles standalone (mirrors the companion section)
\providecommand{\TV}{\mathrm{TV}}
\providecommand{\sgn}{\operatorname{sign}}

The general-capacity section (\S\ref{sec:knowcap-gen}) proved a universal
distribution-free converse and an achievability pinned to two sufficient
regularity classes (Mahalanobis alignment; MLR/log-concave location), leaving the
\emph{complete} characterization --- a single property $\Phi$ with
``computable knowability capacity exists $\iff\Phi$'' --- and any link to causal
structure as open. This section closes the first question at the level of
\emph{existence} with a single property, identifies the exact (small) residue, and
gives a conditional causal correspondence with an honest counterexample.

% ----------------------------------------------------------------------------
\subsection{The problem, abstractly}\label{sec:dich-abstract}
% ----------------------------------------------------------------------------
Fix a measurable space $\mathcal X$, a deployed classifier $g$ and a candidate $f$,
with disagreement region $D=\{x:g(x)\neq f(x)\}$. An \emph{instance} is a pair
$(\vartheta,\eta)$: $\vartheta\in\Theta$ is the \emph{observable parameter} of the
target covariate law $Q_X=Q_\vartheta$ (e.g.\ a shift), and $\eta\in\mathcal H$ is
the \emph{concept} (the target labelling channel $P(Y\mid X)$), which by the
covariate-shift convention does \emph{not} affect $Q_X$. The decision-maker sees an
unlabeled sample $X_{1:n}\sim Q_\vartheta^{\otimes n}$ and must output the sign of
the deployment benefit
\begin{equation}\label{eq:dich-benefit}
  \Delta(\vartheta,\eta)\;=\;R_{Q_\vartheta}(g)-R_{Q_\vartheta}(f)
  \;=\;\int_D \big(1-2\eta(x)\big)\,c_{gf}(x)\,\mathrm dQ_\vartheta(x),
\end{equation}
where $c_{gf}(x)=\pm1$ records which of $g,f$ predicts the larger label on $D$
(a fixed, known function of $g,f$). The admissible set $\mathcal A\subseteq\Theta\times\mathcal H$
encodes prior knowledge (e.g.\ a source calibration plus a mass-drift budget
$\varepsilon$ on the concept). Labels are never observed; a \emph{label-free decision
rule} is any measurable $\hat T:\mathcal X^n\to\{-1,0,+1\}$.

\begin{definition}[Observable reduct]\label{def:dich-obs}
Let $O:\Theta\to\mathcal O$ be the \emph{observable reduct}: $O(\vartheta)$ is a
maximal statistic of $\vartheta$ identifiable from $Q_\vartheta$, i.e.\
$O(\vartheta)=O(\vartheta')\iff Q_\vartheta=Q_{\vartheta'}$. (Concretely $O(\vartheta)$
is the equivalence class $[Q_\vartheta]$; for a location family $O(\vartheta)=\mu$,
for a location--scale family $O=(\mu,\sigma)$.) Because $\eta$ does not affect $Q_X$,
the law of the observed sample depends on $(\vartheta,\eta)$ only through $O(\vartheta)$.
\end{definition}

\begin{definition}[Knowability capacity]\label{def:dich-cap}
A \emph{computable knowability capacity} for $(\{Q_\vartheta\},g,f,\mathcal A)$ is a
pair $(\kappa,\tau)$ with $\kappa:\mathcal O\to\R$ \emph{computable from $Q_X$}
(a measurable function of the observable reduct $O$) and $\tau\in\R$, such that for
every admissible instance the population sign of benefit is label-free identifiable
\emph{exactly} on $\{\kappa(O)>\tau\}$: there is a measurable
$h:\{\kappa>\tau\}\to\{-1,+1\}$ with $\sgn\Delta(\vartheta,\eta)=h(O(\vartheta))$
for all admissible $(\vartheta,\eta)$ with $\kappa(O(\vartheta))>\tau$, and the
label-free minimax error is $\tfrac12$ for every admissible instance with
$\kappa(O(\vartheta))<\tau$. (The finite-$n$ price is the $O(1/\sqrt n)$ Le~Cam
boundary layer of Prop.~\ref{prop:knowcap-finiten} around the level set $\{\kappa=\tau\}$.)
\end{definition}

% ----------------------------------------------------------------------------
\subsection{The single property $\Phi$, and the dichotomy}\label{sec:dich-phi}
% ----------------------------------------------------------------------------
\begin{definition}[The benefit-sign factorization property $\Phi$]\label{def:dich-phi}
We say $(\{Q_\vartheta\},g,f,\mathcal A)$ has property $\Phi$ if the benefit sign
\emph{factors through the observable reduct} on the admissible set:
\begin{equation}\label{eq:dich-phi}
  \boxed{\;
  \begin{aligned}
  \Phi:\quad & O(\vartheta)=O(\vartheta')\ \Longrightarrow\
  \sgn\Delta(\vartheta,\eta)=\sgn\Delta(\vartheta',\eta')\neq 0\\
  & \text{for all admissible }(\vartheta,\eta),(\vartheta',\eta').
  \end{aligned}\;}
\end{equation}
\emph{equivalently}: there is \textbf{no} pair of admissible instances that share the
same observable law $Q_X$ but have opposite (nonzero) benefit sign. Equivalently
again, the \emph{frontier} $\Gamma=\{O:\ \sgn\Delta \text{ changes across admissible
instances with reduct }O\}$ separates $\mathcal O$ into a region of definite $+$ and a
region of definite $-$, and the \emph{frontier margin}
$m(O)=\operatorname{dist}_{\mathrm{mass}}\!\big(\text{admissible band at }O,\ \Gamma\big)$
is well-defined. We say $\Phi$ holds \emph{computably} if, in addition, $m$ is a
computable function of $O$ (Remark~\ref{rmk:dich-compute}).
\end{definition}

The point of \eqref{eq:dich-phi} is that it is a \emph{single} closed condition on
the triple $(\{Q_\vartheta\},g,f)$ and the admissible set --- not a conjunction of a
``unique-flip'' clause and a separate ``monotone-in-nuisance'' clause. Theorem
\ref{thm:dich-main} shows it is exactly the right one.

\begin{theorem}[Knowability dichotomy; converse PROVEN distribution-free, achievability PROVEN under computability]
\label{thm:dich-main}
For any $(\{Q_\vartheta\},g,f,\mathcal A)$ on any measurable space, any dimension, any
family:
\begin{enumerate}\itemsep2pt
\item[\textnormal{(C)}] \textbf{(Converse, PROVEN, distribution-free.)} If $\Phi$
\emph{fails} --- there exist admissible $(\vartheta,\eta),(\vartheta',\eta')$ with
$Q_\vartheta=Q_{\vartheta'}$ and opposite nonzero benefit sign --- then for every $n$
and every label-free $\hat T$,
$\inf_{\hat T}\max_{\{(\vartheta,\eta),(\vartheta',\eta')\}}\Pr[\hat T\text{ wrong}]
   \ge\tfrac12$. No computable capacity exists.
\item[\textnormal{(A)}] \textbf{(Achievability, PROVEN under computability.)} If
$\Phi$ holds \emph{computably}, then $(\kappa,\tau)=(m,\;1)$ \emph{after} rescaling
$m$ by the admissibility budget (so that $\kappa=|m|/\varepsilon$) is a computable
knowability capacity with the \textbf{universal threshold $\tau=1$}: the population
sign of benefit is identifiable exactly on $\{\kappa>1\}=\{$admissible band misses the
frontier$\}$, via the plug-in rule $\hat T=h(\hat O)$, $\hat O$ a consistent estimate
of $O$ from $X_{1:n}$; and the minimax error is $\tfrac12$ on $\{\kappa<1\}$ by part
(C) applied locally.
\end{enumerate}
Consequently: \textbf{a computable knowability capacity exists $\iff$ $\Phi$ holds
computably}, and whenever it exists the capacity object is the frontier margin and
the threshold is the universal constant $1$.
\end{theorem}

\begin{proof}
\emph{(C).} The two worlds generate identical sample laws $Q_\vartheta^{\otimes n}
=Q_{\vartheta'}^{\otimes n}$ (Def.~\ref{def:dich-obs}: $\eta$ does not enter the
observable law and $Q_\vartheta=Q_{\vartheta'}$), so $\TV=0$ for every unlabeled
statistic and every $n$. A label-free sign decision is a test between two simple
hypotheses with opposite correct answers; Le~Cam's two-point bound gives minimax
error $\ge\tfrac12(1-\TV)=\tfrac12$, attained by a constant guess. A capacity would
require a measurable $h(O)$ equal to both $+1$ and $-1$ at the shared reduct, which is
impossible; hence none exists. This is exactly Theorem~\ref{thm:gen-df-conv}, restated
as the necessity half of the dichotomy.

\emph{(A).} Suppose $\Phi$ holds computably. By \eqref{eq:dich-phi} the map
$(\vartheta,\eta)\mapsto\sgn\Delta$ is constant on each fibre $\{O=o\}\cap\mathcal A$
and nonzero there; call the common value $H(o)\in\{-1,+1\}$. Thus $H:\mathcal O\to
\{-1,+1\}$ is well-defined and $\sgn\Delta=H(O)$ on $\mathcal A$. ``Identifiable at
$o$'' for the \emph{full admissible class with reduct $o$} means $H$ is constant on the
\emph{enlarged} fibre obtained by also varying the unobserved $\eta$ within its budget
$\varepsilon$, i.e.\ the admissible concept band at $o$ does not straddle the frontier
$\Gamma$; this is precisely the local event $\{m(o)>\varepsilon\}$, i.e.\
$\kappa(o)=|m(o)|/\varepsilon>1$. On $\{\kappa>1\}$ the rule $\hat T=H(\hat O)$ is
correct in the population and consistent for $\hat O\to O$ (a consistent reduct
estimate exists because $O$ is identifiable from $Q_X$, Def.~\ref{def:dich-obs}); on
$\{\kappa<1\}$ the admissible band straddles $\Gamma$, exhibiting two admissible
instances with the same $O$ (hence same $Q_X$) and opposite sign, so part (C) forces
minimax error $\tfrac12$. The boundary $\{\kappa=1\}$ carries the standard $O(1/\sqrt n)$
Le~Cam layer (Prop.~\ref{prop:knowcap-finiten}) because $\hat O$ has $O(1/\sqrt n)$
fluctuations. Crucially, \emph{no monotonicity of $H$ in $O$ is used}: identifiability
is the local band-vs-frontier event, so the single threshold $\tau=1$ on $\kappa$
governs it even when $H$ flips sign arbitrarily often across $\mathcal O$
(Lemma~\ref{lem:dich-multiflip}).
\end{proof}

\begin{lemma}[$\tau=1$ survives non-monotone, multi-flip frontiers; NUMERICAL]
\label{lem:dich-multiflip}
There are admissible problems in which $\Phi$ holds, the sign map $H(O)$ changes sign
several times across $\mathcal O$ (so the prior section's monotonicity clause R2 is
violated), and yet the single threshold $\tau=1$ on $\kappa=|m(O)|/\varepsilon$
reproduces population identifiability exactly. Hence R1$\wedge$R2 of
\S\ref{sec:knowcap-gen} were \emph{sufficient but not necessary}; the necessary-and-
sufficient content is the single property $\Phi$ (plus computability), with the
margin/threshold the same $\tau=1$.
\emph{(Validated: (B)~a Gaussian-location problem with an observable-but-oscillating
frontier $H$ flipping sign $3$ times --- $0$ mismatches between $\{\kappa>1\}$ and
population identifiability over a $6000$-point grid; and (D)~the strongest stress, a
\emph{two-dimensional} observable reduct $O=(\mu,\sigma)$ with a \emph{curved} frontier
in parameter space --- $0$ mismatches over $\approx\!30{,}000$ random draws and
$330/330$ consistency in the near-boundary shell $0.97<\kappa<1.03$, confirming
identifiability is the \emph{local} band-vs-frontier event governed by the single
threshold $\tau=1$ regardless of the frontier's shape or the reduct's dimension;
\texttt{val\_knowability\_dichotomy.py} Blocks~B,~D.)}
\end{lemma}

\begin{remark}[The exact residue: computability, not regularity]\label{rmk:dich-compute}
The \emph{only} gap between the converse (C) and the achievability (A) is the word
\emph{computably}: (A) needs the frontier margin $m(O)$ to be a computable function of
the observable reduct. This is automatic whenever (i) the shift family is identifiable
from $Q_X$ (so $O$ is recoverable) \emph{and} (ii) the frontier $\Gamma$ is a
measurable function of $O$ obtained by a fixed integral of \eqref{eq:dich-benefit}
(always the case for the location, location--scale, and aligned-Gaussian families, and
for any family where $\Delta$ is a known functional of $Q_\vartheta$). The previously
proposed conditions R1 (unique flip locus) and R2 (monotone-in-nuisance) are
\emph{sufficient} for computability but \emph{not necessary}
(Lemma~\ref{lem:dich-multiflip}). The genuinely open sliver is whether there is an
admissible family for which $\Phi$ holds set-theoretically (the sign factors through
$O$) yet $m(O)$ is \emph{non-computable} (the frontier is a non-measurable / non-
constructive function of $Q_X$); we have no natural example and conjecture none exists
for families with $\Delta$ a fixed integral functional, but we do not prove it
(Conjecture~\ref{conj:dich-compute}).
\end{remark}

\begin{conjecture}[Computability is free for integral-functional families; OPEN]
\label{conj:dich-compute}
If $\Delta(\vartheta,\eta)=\int_D(1-2\eta)c_{gf}\,\mathrm dQ_\vartheta$ with
$\vartheta\mapsto Q_\vartheta$ continuous and the admissible concept budget a
mass-drift ball, then $\Phi\Rightarrow\Phi$-computably; equivalently the frontier
margin $m(O)$ is always a computable functional of $Q_X$. Under this conjecture the
dichotomy is unconditional: \emph{computable capacity exists $\iff\Phi$}. The missing
piece is purely a constructive-measurability statement about the zero level set of a
parametrized integral; it has no probabilistic content.
\end{conjecture}

% ----------------------------------------------------------------------------
\subsection{What $\Phi$ \emph{is}, causally}\label{sec:dich-causal}
% ----------------------------------------------------------------------------
Property $\Phi$ has a transparent causal reading. The benefit's flip locus moves only
when the \emph{unobserved} part of the model --- the concept $\eta=P(Y\mid X)$ ---
moves it; $\Phi$ says exactly that whatever shifts and is benefit-relevant is
\emph{visible in $Q_X$}. We make this precise in a structural causal model (SCM) where
the domain shift is an intervention, and connect it to the causal/anticausal taxonomy
of \citet{scholkopf2012causal} and the information-theoretic UDA analysis of
\citet{wu2024causality}.

\paragraph{SCM and the two directions.} Let the domain shift be an intervention.
\begin{itemize}\itemsep2pt
\item \textbf{Anticausal} ($Y\to X$): the mechanism is $P(X\mid Y)$ and the cause law
is $P(Y)$. By the \emph{independence-of-mechanisms} (ICM) principle
\citep{scholkopf2012causal}, a domain shift perturbs the \emph{cause} marginal $P(Y)$
(label shift) while the mechanism $P(X\mid Y)$ stays invariant. Then $P(Y\mid X)$
\emph{does} change across domains, but it is \emph{recoverable from $Q_X$ alone}
(Bayes with the invariant $P(X\mid Y)$ and the $P(Y)$ identified from the mixture
$Q_X$), \emph{provided the mechanism is informative} (the class-conditionals are
distinguishable, so the mixture is identifiable). Hence the flip locus is a function of
$Q_X$: \textbf{$\Phi$ holds}. This is the population face of \citet{wu2024causality}'s
finding that in anticausal learning the \emph{unlabeled target data dominate}
($O(1/n)$): unlabeled data carry exactly the signal needed to relocate the boundary.
\item \textbf{Causal} ($X\to Y$): the mechanism is $P(Y\mid X)$ \emph{itself}. A
concept-shift intervention perturbs $P(Y\mid X)$ directly and is \emph{invisible} in
$Q_X$ (the cause law may be held fixed). The flip locus then moves with no observable
signal: \textbf{$\Phi$ fails}. This is the population face of
\citet{wu2024causality}'s causal case, where adaptation needs \emph{labelled source}
data and an \emph{unchanged labelling distribution}; if the labelling distribution
changes, unlabeled data are useless --- exactly $\Phi$ failing.
\end{itemize}

\begin{theorem}[Causal correspondence; one inclusion PROVEN, biconditional PROVEN UNDER GENERICITY]
\label{thm:dich-causal}
In the SCM above:
\begin{enumerate}\itemsep2pt
\item[\textnormal{(i)}] \textbf{(Anticausal $\Rightarrow\Phi$, PROVEN under
informativeness.)} If the shift is anticausal with invariant \emph{and informative}
mechanism $P(X\mid Y)$ (the mixture $Q_X$ is identifiable in the cause law), then
$P(Y\mid X)$ is a computable functional of $Q_X$, the flip locus factors through $O$,
and $\Phi$ holds computably; by Theorem~\ref{thm:dich-main} a capacity exists at
$\tau=1$.
\item[\textnormal{(ii)}] \textbf{(Causal-with-mechanism-shift $\Rightarrow\neg\Phi$,
PROVEN under shift-genericity.)} If the shift is causal and \emph{perturbs} $P(Y\mid X)$
within the admissible budget while leaving $Q_X$ fixed, then two admissible instances
share $Q_X$ with opposite benefit sign, so $\Phi$ fails and (Thm~\ref{thm:dich-main}(C))
no capacity exists.
\item[\textnormal{(iii)}] \textbf{(The unconditional biconditional is FALSE.)} There is
a \emph{causal} ($X\to Y$) model with a \emph{stable, observable} mechanism (covariate
shift, no concept drift) for which $\Phi$ \emph{holds}; and an \emph{anticausal} model
with a \emph{degenerate} (uninformative) mechanism for which $\Phi$ \emph{fails}. Hence
``$\Phi\iff$ anticausal'' is not true unconditionally.
\end{enumerate}
Consequently the correspondence is a \textbf{conditional biconditional}:
\[
  \Phi \iff
  \begin{minipage}{0.62\linewidth}\itshape the benefit-relevant shift lies in the
  \emph{observable} factor of the causal factorization, with the unobserved mechanism
  \emph{invariant and informative},\end{minipage}
\]
which \textbf{coincides} with ``anticausal with invariant mechanism'' precisely under
the two genericity assumptions \textnormal{(G1)} a domain shift generically perturbs the
mechanism in the causal direction (so causal $\Rightarrow\neg\Phi$) and \textnormal{(G2)}
the mechanism is generically informative in the anticausal direction (so anticausal
$\Rightarrow\Phi$). Each assumption is \emph{necessary}: dropping (G1) gives the
causal-but-$\Phi$ witness, dropping (G2) the anticausal-but-$\neg\Phi$ witness.
\end{theorem}
\begin{proof}
\emph{(i)} Under invariant informative $P(X\mid Y)$, $Q_X=\sum_y P(Y{=}y)\,P(X\mid Y{=}y)$
is an identifiable mixture, so $P(Y)$ and hence $P(Y\mid X)=P(X\mid Y)P(Y)/q_X$ are
measurable functionals of $Q_X$; the benefit \eqref{eq:dich-benefit} and its flip locus
are therefore functions of $O$, giving $\Phi$ computably and a capacity by
Thm~\ref{thm:dich-main}(A). \emph{(ii)} A within-budget intervention on $P(Y\mid X)$
holding $Q_X$ fixed produces $\eta^+\neq\eta^-$ with the same $Q_X$; choosing the
intervention across the flip locus makes their benefit signs opposite, so
Thm~\ref{thm:dich-main}(C) applies. \emph{(iii)} The two witnesses are exhibited
numerically (Block~C): a Gaussian \emph{covariate}-shift model ($X\to Y$, $\eta$ fixed)
where $\sgn\Delta$ is a deterministic function of the observable shift $\mu$
($0$ capacity mismatches over $6{,}000$ draws), and an anticausal model with coincident
class-conditionals ($P(X\mid Y)$ independent of $Y$) where $Q_X$ is invariant to the
prior $\pi$ yet $\sgn\Delta=\sgn(1-2\pi)$ takes both signs.
\emph{(Validated: C1 anticausal+invariant $0/6000$ mismatches; C2 causal+shifted both
signs at $\TV{=}0$; C3 causal+stable $0/6000$ mismatches [causal-but-$\Phi$]; C4
anticausal+degenerate both signs at one $Q_X$ [anticausal-but-$\neg\Phi$];
\texttt{val\_knowability\_dichotomy.py} Block~C.)}
\end{proof}

\begin{remark}[What is new vs.\ \citealp{scholkopf2012causal,wu2024causality}]
\label{rmk:dich-novelty}
\citet{scholkopf2012causal} introduced the causal/anticausal split and the ICM
principle, arguing \emph{qualitatively} that anticausal problems admit semi-/un-
supervised gains while causal ones do not. \citet{wu2024causality} made this
\emph{quantitative} for \emph{rates}: with labelled source data the excess risk decays
as $O(1/m)$ (causal, invariant labelling) versus $O(1/n)$ (anticausal, unlabeled-driven).
Both presuppose labels (in the source) and ask \emph{how fast} one learns. The present
result is on a \emph{different axis}: it is \textbf{fully label-free} (no source labels)
and asks \emph{whether the sign of a deployment benefit is knowable at all}, answering
with an \emph{exact threshold} $\tau=1$ on a computable capacity. The contribution is
(a) the single structural property $\Phi$ (benefit-sign factorization) and the
\emph{distribution-free} existence dichotomy (Thm~\ref{thm:dich-main}), which has no
analogue in either prior work; and (b) the identification of $\Phi$ with an
\emph{observable-shift / invariant-informative-mechanism} condition that recovers the
causal/anticausal split \emph{only under two explicit genericities}
(Thm~\ref{thm:dich-causal}), together with the two counterexamples showing the
unconditional ``knowable $\iff$ anticausal'' slogan is false. To our knowledge the
label-free knowability dichotomy and its conditional --- not unconditional --- causal
map are new.
\end{remark}

\begin{remark}[Honest scope and the wall]\label{rmk:dich-scope}
\textbf{Q1 (single $\Phi$).} \emph{Proven:} a single property $\Phi$ (benefit-sign
factorization through the observable reduct) yields a \emph{complete, distribution-free}
dichotomy for the \emph{existence} of a label-free knowability capacity, with the
\emph{universal} threshold $\tau=1$; the converse is the universal Le~Cam bound and the
achievability holds whenever $\Phi$ is \emph{computable}. We further prove $\tau=1$
survives non-monotone, multi-flip frontiers, so the prior section's R1$\wedge$R2 were
\emph{sufficient but not necessary}: the true characterization is the single $\Phi$, not
a conjunction. \emph{Not proven:} that $\Phi$ (set-theoretic factorization) always
\emph{implies} $\Phi$-computability --- i.e.\ that the frontier margin is always a
computable functional of $Q_X$. This residue (Conjecture~\ref{conj:dich-compute}) is
purely a constructive-measurability question with \emph{no probabilistic content}, and
is automatically discharged for every concrete family we treat (location,
location--scale, aligned-Gaussian, Gaussian mixture). So at the level of \emph{existence}
the dichotomy is single-property and complete; only if one demands the capacity be a
\emph{closed-form} computable margin does a short second clause (computability) appear,
and we have no example where it genuinely bites.

\textbf{Q2 (causal map).} \emph{Proven:} the two natural inclusions (anticausal+invariant
+informative $\Rightarrow\Phi$; causal+mechanism-shift $\Rightarrow\neg\Phi$) and that
the \emph{unconditional} biconditional is \emph{false} (two explicit counterexamples).
\emph{Proven only under genericity:} the biconditional ``$\Phi\iff$ anticausal-with-
invariant-mechanism'', valid under (G1) mechanism-shift-genericity and (G2) mechanism-
informativeness, each shown necessary. \emph{The wall:} we cannot remove (G1)/(G2);
$\Phi$ is genuinely a statement about \emph{which factor of the distribution is
observable and shifting}, and the causal \emph{direction} controls that only
generically. A non-degenerate-shift / faithfulness assumption is exactly the bridge, and
formalizing the weakest such assumption that makes the biconditional exact (rather than
generic) is left open.

\textbf{Verdict.} This is a \emph{strong, fundamental, and honestly-bounded} step, not an
unconditional paradigm theorem: Q1 reaches a genuine single-property distribution-free
dichotomy for existence (with a residue that is non-probabilistic and never observed to
bite), and Q2 reaches a \emph{conditional} causal correspondence with the unconditional
slogan explicitly disproven. The fundamental object --- one property, one universal
threshold $\tau=1$, mapped to the observability of the shifting causal factor --- is, we
believe, the right one; the remaining gap is the genericity in the causal map and the
constructive-measurability sliver in Q1, both precisely located.
\end{remark}
